%%%%%%%%%%%%%%%%%%%%%%%%%%%%%%%%%%%%%%%%%
% Medium Length Professional CV
% LaTeX Template
% Version 2.0 (8/5/13)
%
% This template has been downloaded from:
% http://www.LaTeXTemplates.com
%
% Original author:
% Rishi Shah 
%
% Important note:
% This template requires the resume.cls file to be in the same directory as the
% .tex file. The resume.cls file provides the resume style used for structuring the
% document.
%
%%%%%%%%%%%%%%%%%%%%%%%%%%%%%%%%%%%%%%%%%

% TODO: Update Mphil supervision
% TODO: add Schmidt IX interview thing

%----------------------------------------------------------------------------------------
%	PACKAGES AND OTHER DOCUMENT CONFIGURATIONS
%----------------------------------------------------------------------------------------

\documentclass{resume} % Use the custom resume.cls style

\usepackage[left=0.75in,top=0.4in,right=0.75in,bottom=0.4in]{geometry} % Document margins
\usepackage{ragged2e}
\usepackage[
colorlinks=true,
linkcolor={green},
urlcolor=blue,
]{hyperref}
\usepackage{makecell}
\newcommand{\tab}[1]{\hspace{.2667\textwidth}\rlap{#1}}
\newcommand{\itab}[1]{\hspace{0em}\rlap{#1}}
\name{Gergely Flamich} % Your name
\address{Imperial College London, London, UK} % Your address
%\address{123 Pleasant Lane \\ City, State 12345} % Your secondary addess (optional)
\address{+44\,74\,62\,957\,131\, \\ \,g.flamich@imperial.ac.uk\, \\ \,\url{https://gergely-flamich.github.io}} % Your phone number and email

\begin{document}

% ----------------------------------------------------------------------------------------
%	PERSONAL STATEMENT
% ----------------------------------------------------------------------------------------

 \begin{rSection}{Personal Statement}
 Since August 2025, I have been an
 \href{https://www.imperial.ac.uk/research-and-innovation/research-office/funder-information/research-fellowships/irf/}{Imperial
   College Research Fellow} at Imperial College London.
 I have deep expertise in machine learning with strong practical and
 theoretical research and coding skills.
 My current work focuses on \textbf{neural data compression} using generative
 models such as diffusion models and variational autoencoders, and related
 \textbf{information theory}, mainly \textbf{relative entropy coding}.

\end{rSection}
%----------------------------------------------------------------------------------------
%	PROFESSIONAL EXPERIENCE SECTION
%----------------------------------------------------------------------------------------



%----------------------------------------------------------------------------------------
%	WORK EXPERIENCE SECTION
%----------------------------------------------------------------------------------------

\begin{rSection}{Professional Experience}
\begin{rSubsection}{Imperial College Research Fellow (Research Associate until
    July 2025)}{March 2025 - Now}{Host: Dr Deniz G{\"u}nd{\"u}z}{\textit{Imperial College London}}
\item Characterised the absolute best theoretically achievable data compression using a new fundamental measure of information called the
  \textit{functional information} (a generalisation of entropy). Paper in
  preparation for IEEE Transactions on Information Theory.
\item Developed the first data compression algorithm that has \textbf{provably
    optimal compression performance}. Implementation and analysis mainly in \texttt{Jax}. Paper to appear at IEEE ITW 2026.
\item Implemented scalable, privacy-preserving image compression algorithm using diffusion
  models that compresses data \textbf{$\times$20-40 more} than competing approaches.
  Implementation mostly in \texttt{PyTorch}.
  Workshop paper to appear at ICML 2026, full paper in preparation for ICLR 2027.
\item Co-organised the \href{https://learn-to-compress-workshop-isit.github.io}{Learn to Compress \& Compress to Learn
Workshop} at the International Symposium on Information Theory (ISIT) 2025;
currently assisting in organisation of the 2026 rendition.
\item Was elected as \href{https://www.imperial.ac.uk/research-and-innovation/research-office/funder-information/research-fellowships/irf/our-cohorts/}{one of 16 new research fellows at Imperial College} in August 2025, based
  on an independently written research proposal.
  As the \href{https://www.imperial.ac.uk/research-and-innovation/research-office/funder-information/research-fellowships/irf/}{University's website states}, the fellowship `is designed for exceptional early career researchers who are prepared to step into independence.'
  The fellowship covers salary as well as £30,000 for equipment and travel expenses.
\item Since start of postdoc, 3 papers to appear at information theory and ML
  venues, of which one received spotlight award; 3
  submitted for peer review.
  \href{https://www.imperial.ac.uk/early-career-researcher-institute/learning-and-development/cornerstone/supervisors-guidebook/cpd/assistant-supervisors/}{Assistant supervisor} to first-year PhD student ({\"O}. G{\"u}ner) and
  supervising an MPhil research project at Cambridge on video
  compression with diffusion models.
\end{rSubsection}

\begin{rSubsection}{Student Researcher: Neural Machine Translation}{Dec
      2024 - Feb 2025}{Hosts: Dr David Vilar, Dr Jan-Thorsten
    Peter and Dr Markus Freitag}{\textit{Google Berlin}}
\item Implemented discrete diffusion models in \texttt{Jax} and \texttt{T5X} to speed up neural machine translation.
\item Developed theory of the accuracy-naturalness tradeoff in translation and
 demonstrated that current methodology for automatic evaluation of
  translation systems is inadequate and needs to be updated.
\item Work resulted in a paper published and presented at
  the \href{https://openreview.net/forum?id=d9EkgbZZH9}{Conference on Language Modeling 2025}.
\end{rSubsection}

\begin{rSubsection}{Student Researcher: Relative Entropy Coding for Practical Data Compression}{July
      2022 - Dec 2022}{Host: Dr Lucas Theis}{\textit{Google Brain, London}}
\item Developed data compression algorithm that is significantly
  faster than previous comparable methods. Implementation in \texttt{TensorFlow}.
\item Published and presented work at
  \href{https://ieeexplore.ieee.org/document/10206725}{IEEE International
    Symposium on Information Theory 2023}.
\item Contributed the noncentral $\chi^2$ distribution to
  the open-source \href{https://www.tensorflow.org/probability}{TensorFlow Probability} package.
\end{rSubsection}

\begin{rSubsection}{PhD in Advanced Machine Learning}{Oct 2020 - July 2025}{Supervisor: Dr Jos\'e Miguel Hern\'andez-Lobato}{\textit{University of Cambridge}}
\item Developed new theoretical foundations for lossy data compression using
  randomised codes.
\item Applied the theoretical findings to develop lossy data compression
  algorithms using generative models, such as variational autoencoders and
  Bayesian implicit neural representations.
\item Published 14 papers (first author on 8) with 11 coauthors at top AI/ML and information theory
  conferences, such as NeurIPS, ICLR, ICML, COLM and ISIT.
  Received \textbf{spotlight award (top 10\% accepted papers)} at NeurIPS for work on
  efficient data compression with Bayesian INRs.
  Proposed and supervised three MPhil research projects and two undergraduate
  research projects.
\end{rSubsection}

%\vspace{-0.7cm}
\begin{rSubsection}{Research Assistant}{Oct 2019 - July 2020}{Supervisor: Dr Jos\'e Miguel Hern\'andez-Lobato}{University of Cambridge}
\item Implemented and optimized black-box optimization pipeline to find optimal
  hardware designs in collaboration with researchers at Harvard and Intel.
  Implementation in \texttt{TensorFlow} and \texttt{GPFlow}.
\item Our method consistently found better configurations faster than standard
  methods used in industry.
\end{rSubsection}

\end{rSection}

%----------------------------------------------------------------------------------------
%	EDUCATION SECTION
%----------------------------------------------------------------------------------------

\begin{rSection}{Education}

\rSubsectionNoContent{PhD in Advanced Machine Learning}{Oct 2020 - July
    2025}{\textbf{Supervisor:} Dr Jos\'e Miguel Hern\'andez-Lobato}{University
    of Cambridge}

\rSubsectionNoContent{\bf MPhil in Machine Learning and Machine Intelligence}{Oct
    2018 - Oct 2019}{Graduated with \textbf{Commendation} (roughly top 20\% of class)}{University
    of Cambridge}

\rSubsectionNoContent{\bf BSc Joint Honours in Mathematics and Computer Science}{Oct
    2014 - Jun 2018}{\textit{Graduated as \textbf{Valedictorian} in Computer Science, with First Class Honours}}{University
    of St Andrews}
\end{rSection}

%----------------------------------------------------------------------------------------
%	Academic Achievements
%----------------------------------------------------------------------------------------

\begin{rSection}{Academic Achievements}
  \begin{tabular}{rll}
    2026     & Spotlight Paper Award & \href{https://learn-to-compress-workshop-isit.github.io/2026/schedule/}{Learn to Compress Workshop @ ISIT 2026} \\[1ex]
    2025     & \makecell[l]{\href{https://www.imperial.ac.uk/research-and-innovation/research-office/funder-information/research-fellowships/irf/}{Imperial College} \\
    \href{https://www.imperial.ac.uk/research-and-innovation/research-office/funder-information/research-fellowships/irf/}{Research Fellowship}} & \makecell[l]{Equivalent to Leverhulme Trust \& Wellcome Trust Early Career \\ Fellowships. Success rate: 17.98\% of applicants (2025).} \\[2ex]
%    2025     & Expert Reviewer & \href{https://jmlr.org/tmlr/expert-reviewers.html}{Transactions on Machine Learning Research} \\
    2025     & \makecell[l]{Runner-up for the Jack Keil Wolf \\ student paper award} & ISIT 2025 \\[2ex]
%    2024     & Expert Reviewer & \href{https://jmlr.org/tmlr/expert-reviewers.html}{Transactions on Machine Learning Research} \\
    2024     & Spotlight Paper Award & \href{https://learn-to-compress-workshop-isit.github.io/2024/schedule/}{Learn to Compress Workshop @ ISIT 2024} \\[1ex]
    2023     & Spotlight Paper Award & \href{https://openreview.net/forum?id=5otj6QKUMI&noteId=tVmGxIMDWn}{NeurIPS 2023}, a top venue for machine learning research. \\[1ex]
    2021     & Finalist & \href{https://www.qualcomm.com/research/university-relations/innovation-fellowship/finalists}{Qualcomm Innovation Fellowship}
  \end{tabular}
\end{rSection}

\begin{rSection}{Technical Skills}

\begin{tabular}{ @{} >{\bfseries}l @{\hspace{6ex}} l }
Programming Languages \ & Python, Javascript, Java, Haskell, Matlab, C, C++, \LaTeX \\
ML Frameworks \& APIs & Jax AI Stack, TensorFlow, TensorFlow Probability, SciPy \\
\end{tabular}

\end{rSection}


%----------------------------------------------------------------------------------------
% Publications
%----------------------------------------------------------------------------------------

\begin{rSection}{Selected Publications}\setlength\itemsep{-0.5em} \itemsep -6pt
% \item \textbullet\,\, \textbf{G. Flamich}, {\"O}. G{\"u}ner, Y. Liu and D.
%   G{\"u}nd{\"u}z (2026). Scalable Differentially Private Data Compression via
% Diffusion and Stochastic Codes. To appear at \textit{ICML 2026 Workshop on
% Structured Probabilistic Inference \& Generative Modeling}.
\item \textbullet\,\,\textbf{G. Flamich} and S. Hill (2026).
  Singular relative entropy coding with bits-back rejection sampling. To appear
  in \textit{Learn to Compress \& Compress to Learn Workshop @ ISIT 2026}.
  Received \textbf{Spotlight award}.
% \item \textbullet\,\,\textbf{G. Flamich} and D. G{\"u}nd{\"u}z (2026).
%   Data Compression with Stochastic Codes. To appear in the \textit{IEEE BITS Magazine (2026)}.
\item \textbullet\,\,\textbf{G. Flamich}, S. M. Sriramu and A. B. Wagner (2025). The
  Redundancy of Non-Singular Channel Simulation. In \textit{ISIT 2025}.
  Runner-up to the Jack Keil Wolf student paper award (\textbf{top 6 student paper})
% \item \textbullet\,\,\textbf{G. Flamich} and L. Wells. Some Notes on the Sample
%   Complexity of Approximate Channel Simulation. In \textit{First
%     ’Learn to Compress’ Workshop@ ISIT 2024}. Received \textbf{Spotlight award}
%   (top 4 of 16 accepted works).
\item \textbullet\,\,J. He$^\dagger$, \textbf{G. Flamich}$^\dagger$, Z. Guo, J. M. Hern\'andez-Lobato.
    RECOMBINER: Robust and Enhanced Compression with Bayesian Implicit Neural
    Representations. In \textit{ICLR 2024}.
\item \textbullet\,\,\textbf{G. Flamich}.
    Greedy Poisson Rejection Sampling. In \textit{NeurIPS 2023}.
  \item \textbullet\,\,\textbf{G. Flamich}$^\dagger$, Z. Guo$^\dagger$, J. He, Z. Chen, J. M. Hern\'andez-Lobato.
  Compression with Bayesian Implicit Neural Representations. In \textit{NeurIPS
    2023}. Received \textbf{Spotlight award} (top 10\% of accepted papers).
\item \textbullet\,\,\textbf{G. Flamich}$^\dagger$, S. Markou$^\dagger$, J. M. Hern\'andez-Lobato.
    Fast Relative Entropy Coding with A* coding. In \textit{ICML 2022}.
  \item \textbullet\,\,\textbf{G. Flamich}$^\dagger$, M. Havasi$^\dagger$, J. M. Hern\'andez-Lobato.
    Compressing Images by Encoding Their Latent Representations with Relative
    Entropy Coding. In \textit{NeurIPS 2020}.
    \\
    $^\dagger$ equal contribution.
\end{rSection}


\end{document}
